% =====================================================================
\section{How AI changes the classical side}
\label{sec:ai}

The classical responses of Section~\ref{sec:devices} came from tensor networks, operator-propagation and variational methods. It does not follow that AI contributes only by imitation. Four roles appear in the evidence, and merging them is the source of much confusion about whether ``AI'' competes with quantum computing.

\subsection{Direct solvers}
A neural quantum state is a neural network used as a trial wavefunction and optimized against the Schr\"odinger equation by variational Monte Carlo~\cite{carleo2017,carleo2019}. In the examples considered here the network parameters are optimized against the Hamiltonian through a variational energy objective, with no labelled solutions; other neural-state workflows add pretraining or data. A flexible ansatz still has limits of expressivity, sampling and optimization. On the frustrated $J_1$--$J_2$ Heisenberg model at $J_2/J_1=0.5$ on a $10\times10$ lattice, a transformer wavefunction reaches an energy per site of $-0.497634(1)$, below a DMRG value of $-0.495530$, in four days on twenty GPUs~\cite{rende2024}; that is a measured optimization run, not the cost of reaching a certified accuracy. Transformer wavefunctions have been applied to doped Hubbard models with longer-range hopping~\cite{gu2026}, neural variational Monte Carlo to molecular excited states~\cite{pfau2024,schatzle2026}, and hidden-nucleon networks to nuclei up to mass number 20, giving binding energies, charge radii and magnetic moments for interactions from pionless effective field theory and related models~\cite{lovato2022,gnech2024,lovato2026}. For spin glasses, autoregressive networks annealed variationally outperform both simulated annealing and a classical emulation of quantum annealing over the sizes studied~\cite{hibat2021,hibat2021b}.

These are the classical methods a fair comparison must now include, and they come with their own limits. A Monte Carlo error bar is not a bound on the bias of the ansatz, and a lower variational energy is not a certified interval. Results depend on optimizer and random seed; long-time dynamics is unstable on current evidence~\cite{tindall2025}; and the cost of reaching a fixed, verified accuracy is rarely reported. None of the results above is on the instance of a quantum estimate reviewed in Section~\ref{sec:ft}: the lattice results use other geometries and the molecular results other Hamiltonians. They are candidates, not matches.

\subsection{Amortization}
A model trained once can answer many queries, and this changes the comparison whenever the workload is large. Fig.~\ref{fig:funnel} shows a complete example. A Microsoft--PNNL workflow combined learned models, density functional theory and molecular dynamics to reduce 32.6 million candidate solid electrolytes to 18 in under 80 hours, and the leading family was then synthesized and its ionic conductivity measured~\cite{chen2024}. Graph networks have likewise predicted 2.2 million crystal structures stable against earlier data, about 381{,}000 of them on the updated convex hull~\cite{merchant2023}; predicted stability is not synthesizability.

\begin{figure}[htbp]
\centering
\begin{tikzpicture}
\begin{axis}[
  xbar, xmode=log, log origin=infty, width=0.72\textwidth, height=6.0cm,
  xmin=5, xmax=3e8, y dir=reverse, ytick={1,...,6},
  yticklabels={Generated (AI), Predicted stable (ML), Functional filter (ML), DFT + ML-MD (HPC), Practicality filter, Predicted new candidates},
  yticklabel style={font=\footnotesize}, xlabel={\footnotesize Candidates remaining (log scale)},
  xticklabel style={font=\footnotesize}, bar width=10pt, enlarge y limits=0.1,
  nodes near coords, point meta=rawx, every node near coord/.append style={font=\scriptsize, anchor=west},
  axis x line*=bottom, axis y line*=left, xmajorgrids, grid style={black!10}]
\addplot[fill=qcN, draw=none, bar shift=0pt] table[x=x,y=y] {figdata/funnel_screen.dat};
\addplot[fill=qcF, draw=qcQ, line width=0.3pt, bar shift=0pt, postaction={pattern=north east lines, pattern color=white}] table[x=x,y=y] {figdata/funnel_new.dat};
\node[font=\scriptsize\itshape, anchor=west, align=left, black!80] at (axis cs:300,5.6) {after the screen (not in the 80 h):\\one family synthesized, structure\\and ionic conductivity measured};
\end{axis}
\end{tikzpicture}
\caption[An AI+HPC workflow: solid-electrolyte screening]{A complete classical workflow: solid-electrolyte screening by machine-learned models, density functional theory and molecular dynamics on about 1{,}000 cloud virtual machines, from 32.6 million candidates to 18 predicted new compositions in under 80 hours of wall time~\cite{chen2024}. Plain bars are screening stages; the last bar, hatched, is the set of predicted new compositions. The step from 23 to 18 is a novelty filter against known compounds, not a validation rate. Synthesis and measurement of the Na$_x$Li$_{3-x}$YCl$_6$ series followed the screen and are a separate endpoint. No quantum processor was used. Stage counts are in \texttt{values.csv}.}
\label{fig:funnel}
\end{figure}


These workflows succeed without solving the strongly correlated problems proposed for quantum computers, so they are not solver comparisons. Their proper measure is the quality, cost and reliability of the decision they support. A quantum solver could matter at the few decision points where the error of density functional theory dominates the error budget, and the test is to identify such a point. Theory supports a narrower statement: given measurement data from related systems, classical learning predicts ground-state properties efficiently within a gapped phase~\cite{huang2022}. The data must come from somewhere, possibly a quantum device, and with a quantum memory there is an exponential advantage in the number of experiments needed for certain learning tasks~\cite{huang2022exp}. Amortization therefore cuts both ways. Quantum results can also be reused, as training labels or across observables.

\subsection{Reformulation}
A learned model can replace the hard computation with a different one that answers the same scientific question. The Skala functional learns the exchange--correlation term of density functional theory from high-accuracy wavefunction data and reaches chemical accuracy on a standard set of atomization energies, with weighted errors of 3.89 and 2.72\,kcal/mol on a broad thermochemistry benchmark in its two releases, at the cost of semi-local density functional theory~\cite{luise2025,skala10,skala11}. AlphaFold learns the map from sequence to protein structure~\cite{jumper2021}. A reformulation is only as good as its validated domain. Skala is validated on main-group molecules like those in its training data and gives no per-prediction uncertainty. For transition-metal and multireference chemistry the labels themselves are the bottleneck, and fitting labels that share an unknown systematic bias does not remove the bias. Accurate labels for strongly correlated systems are one place where a quantum computer could feed a classical workflow, if their value to the downstream decision can be shown. A failure of a learned model outside its domain creates a need for a better solver, not automatically a quantum one.

\subsection{Theory}
In November 2026 a proof was posted of a 2014 conjecture of Aaronson and Ambainis in query complexity, with the central inequality checked in the Lean proof assistant~\cite{aa2014,turturean2026}. By its author's account, AI models produced the mathematical statements, the first draft and the formalization.

\begin{rulebox}{A query-complexity result, and what it does not say}
A query model charges for looking up input data and lets all other computation run free. The claimed proof limits one kind of quantum query advantage: for yes-or-no questions, a classical procedure can approximate the acceptance probability of any quantum procedure on most uniformly random inputs using polynomially many queries. This extends a relation long known for the worst case~\cite{beals2001,abkrt2021}. It does not bound the classical procedure's running time, it does not cover sampling or problems with a promise such as period finding, and it does not rank any chemistry, dynamics or sampling workload in this survey. Its relevance here is methodological: the promise on the input and the model of data access decide what a speedup claim means. The proof is a preprint; the formal statement, its verification status and a numerical illustration are in Appendix~\ref{app:query}.
\end{rulebox}

We record this as mathematics about the boundary between quantum and classical computing, produced with AI assistance. It is not a case of AI+HPC against quantum computing, and it is not evidence about the scientific problems reviewed here.

\subsection{AI methods need baselines too}
The discipline asked of quantum claims applies equally to learned methods. A physics-inspired graph neural network was reported to match or beat existing solvers for maximum independent set at million-variable scale; a simple greedy algorithm then found larger sets and ran $10^4$ times faster at $10^6$ nodes, and the network's authors replied that this sparse problem is unrepresentative~\cite{angelini2022,schuetz2023}. One weak neural baseline does not make neural optimization uncompetitive, any more than one weak classical baseline establishes a quantum advantage. Learned-model results in this survey are therefore read with the same questions as quantum ones: what domain was validated, what the labels cost, how uncertainty is calibrated, and what the strongest non-learned baseline achieves.
